% GENERATED FILE. Edit canonical lessons, then run npm run generate:books.
% canonical-source-hash: fnv1a64:47df78dcc765e226
% canonical-lessons: UR-C163-fuzul, UR-C163-majbur, UR-C163-azad, UR-C163-pak, UR-C163-sharminda

\chapter{Pointless, Forced, Free, Pure, Embarrassed}
\label{ch:ur-pointless-forced-free-pure-embarrassed}
\hlchaptermodality{163}

\begin{chapteropening}
\textbf{By the end of this chapter:} \emph{I can say pointless, forced, free, pure and embarrassed.}

Complete the last lesson of chapter 163: I can say pointless, forced, free, pure and embarrassed.
\end{chapteropening}

\section[\texorpdfstring{\ur{فضول} (fuzūl)}{فضول (fuzūl)}]{\ur{فضول} (fuzūl) — pointless, wasteful}

\label{lesson:UR-C163-fuzul}



\begin{quote}
Before the new one: say the Urdu for Indian, then the Urdu for grey.
\end{quote}

\subsection*{You'll want to know: \ur{فضول}}
\textbf{\ur{فضول}} — \emph{fuzūl} — \textquotedblleft{}pointless, wasteful\textquotedblright{}.

\begin{cousinweb}
\textbf{\ur{فضول}} (\emph{fuzūl}) is a loan from Arabic \emph{fuḍūl}, \textquotedblleft{}excesses, what is left over\textquotedblright{}, which reached Urdu through Persian. Its Arabic root is \emph{f-ḍ-l} (Lane). It is the plural of \emph{faḍl}, \textquotedblleft{}excess\textquotedblright{}.
\end{cousinweb}

\begin{grammarlens}[title={What you've built}]
One more word for this chapter.
\end{grammarlens}

\subsection*{Your turn}
\begin{itemize}
\raggedright
  \item \emph{Say it:} \emph{fuzūl}
  \item \emph{Say it:} \emph{fuzūl}, once more
  \item \emph{Say it:} say \emph{fuzūl}
  \item \emph{Recall:} say the Urdu for Indian, then the Urdu for grey, then say \emph{fuzūl} again
\end{itemize}

\begin{tcolorbox}[breakable,colback=teal!4,colframe=teal!35!black,title={Before you move on}]
What does \ur{فضول} mean? (\textquotedblleft{}Pointless, wasteful\textquotedblright{}.) Say it once more.
\end{tcolorbox}
\section[\texorpdfstring{\ur{مجبور} (majbūr)}{مجبور (majbūr)}]{\ur{مجبور} (majbūr) — forced, helpless}

\label{lesson:UR-C163-majbur}



\begin{quote}
Before the new one: say the Urdu for grey, then the Urdu for pointless.
\end{quote}

\subsection*{You'll want to know: \ur{مجبور}}
\textbf{\ur{مجبور}} — \emph{majbūr} — \textquotedblleft{}forced, helpless\textquotedblright{}.

\begin{cousinweb}
\textbf{\ur{مجبور}} (\emph{majbūr}) is a loan from Arabic \emph{majbūr}, \textquotedblleft{}compelled\textquotedblright{}, which reached Urdu through Persian. Its Arabic root is \emph{j-b-r} (Lane).
\end{cousinweb}

\begin{grammarlens}[title={What you've built}]
One more word for this chapter.
\end{grammarlens}

\subsection*{Your turn}
\begin{itemize}
\raggedright
  \item \emph{Say it:} \emph{majbūr}
  \item \emph{Say it:} \emph{majbūr}, once more
  \item \emph{Say it:} say \emph{majbūr}
  \item \emph{Recall:} say the Urdu for grey, then the Urdu for pointless, then say \emph{majbūr} again
\end{itemize}

\begin{tcolorbox}[breakable,colback=teal!4,colframe=teal!35!black,title={Before you move on}]
What does \ur{مجبور} mean? (\textquotedblleft{}Forced, helpless\textquotedblright{}.) Say it once more.
\end{tcolorbox}
\section[\texorpdfstring{\ur{آزاد} (āzād)}{آزاد (āzād)}]{\ur{آزاد} (āzād) — free (not captive)}

\label{lesson:UR-C163-azad}



\begin{quote}
Before the new one: say the Urdu for pointless, then the Urdu for forced.
\end{quote}

\subsection*{You'll want to know: \ur{آزاد}}
\textbf{\ur{آزاد}} — \emph{āzād} — \textquotedblleft{}free (not captive)\textquotedblright{}.

\begin{cousinweb}
\textbf{\ur{آزاد}} (\emph{āzād}) is a loan, usually traced to Persian \emph{āzād}, \textquotedblleft{}free\textquotedblright{}.
\end{cousinweb}

\begin{grammarlens}[title={What you've built}]
One more word for this chapter.
\end{grammarlens}

\subsection*{Your turn}
\begin{itemize}
\raggedright
  \item \emph{Say it:} \emph{āzād}
  \item \emph{Say it:} \emph{āzād}, once more
  \item \emph{Say it:} say \emph{āzād}
  \item \emph{Recall:} say the Urdu for pointless, then the Urdu for forced, then say \emph{āzād} again
\end{itemize}

\begin{tcolorbox}[breakable,colback=teal!4,colframe=teal!35!black,title={Before you move on}]
What does \ur{آزاد} mean? (\textquotedblleft{}Free (not captive)\textquotedblright{}.) Say it once more.
\end{tcolorbox}
\section[\texorpdfstring{\ur{پاک} (pāk)}{پاک (pāk)}]{\ur{پاک} (pāk) — pure, clean}

\label{lesson:UR-C163-pak}



\begin{quote}
Before the new one: say the Urdu for forced, then the Urdu for free.
\end{quote}

\subsection*{You'll want to know: \ur{پاک}}
\textbf{\ur{پاک}} — \emph{pāk} — \textquotedblleft{}pure, clean\textquotedblright{}.

\begin{cousinweb}
\textbf{\ur{پاک}} (\emph{pāk}) is a loan, usually traced to Persian \emph{pāk}, \textquotedblleft{}pure, clean\textquotedblright{}.
\end{cousinweb}

\begin{grammarlens}[title={What you've built}]
One more word for this chapter.
\end{grammarlens}

\subsection*{Your turn}
\begin{itemize}
\raggedright
  \item \emph{Say it:} \emph{pāk}
  \item \emph{Say it:} \emph{pāk}, once more
  \item \emph{Say it:} say \emph{pāk}
  \item \emph{Recall:} say the Urdu for forced, then the Urdu for free, then say \emph{pāk} again
\end{itemize}

\begin{tcolorbox}[breakable,colback=teal!4,colframe=teal!35!black,title={Before you move on}]
What does \ur{پاک} mean? (\textquotedblleft{}Pure, clean\textquotedblright{}.) Say it once more.
\end{tcolorbox}
\section[\texorpdfstring{\ur{شرمندہ} (sharmindā)}{شرمندہ (sharmindā)}]{\ur{شرمندہ} (sharmindā) — embarrassed}

\label{lesson:UR-C163-sharminda}



\begin{quote}
Before the new one: say the Urdu for free, then the Urdu for pure.
\end{quote}

\subsection*{You'll want to know: \ur{شرمندہ}}
\textbf{\ur{شرمندہ}} — \emph{sharmindā} — \textquotedblleft{}embarrassed\textquotedblright{}.

\begin{cousinweb}
\textbf{\ur{شرمندہ}} (\emph{sharmindā}) is a loan, usually traced to Persian \emph{sharminda}, \textquotedblleft{}ashamed\textquotedblright{}. It is built on \emph{sharm}, \textquotedblleft{}shame\textquotedblright{}.
\end{cousinweb}

\begin{grammarlens}[title={What you've built}]
One more word for this chapter.
\end{grammarlens}

\subsection*{Your turn}
\begin{itemize}
\raggedright
  \item \emph{Say it:} \emph{sharmindā}
  \item \emph{Say it:} \emph{sharmindā}, once more
  \item \emph{Say it:} say \emph{sharmindā}
  \item \emph{Recall:} say the Urdu for free, then the Urdu for pure, then say \emph{sharmindā} again
\end{itemize}

\begin{tcolorbox}[breakable,colback=teal!4,colframe=teal!35!black,title={Before you move on}]
What does \ur{شرمندہ} mean? (\textquotedblleft{}Embarrassed\textquotedblright{}.) Say it once more.
\end{tcolorbox}
